\documentclass[11pt]{article}

% =========================================================
% PACKAGES
% =========================================================
\usepackage[margin=1in]{geometry}
\usepackage{amsmath,amssymb,amsfonts,amsthm}
\usepackage{graphicx}
\usepackage{float}
\usepackage{subcaption}
\usepackage{caption}
\usepackage{booktabs}
\usepackage{longtable}
\usepackage{multirow}
\usepackage{array}
\usepackage{xcolor}
\usepackage{hyperref}
\usepackage{enumitem}
\usepackage{parskip}

% =========================================================
% FORMATTING
% =========================================================
\hypersetup{
    colorlinks=true,
    linkcolor=blue,
    citecolor=blue,
    urlcolor=blue
}

\setlist[itemize]{leftmargin=1.5em}
\setlist[enumerate]{leftmargin=1.5em}

\graphicspath{
    {../results/figures/}
    {../../results/figures/01_dgp_eda/}
    {../../results/figures/02_var_deepar_fit/}
    {../../results/figures/03_residuals/}
    {../../results/figures/04_innovations/}
    {../../results/figures/05_diffusion/}
    {../../results/figures/06_forecasts/}
    {../../results/figures/07_calibration_evaluation/}
}

% =========================================================
% TITLE
% =========================================================
\title{
Flexible Innovation Modeling for Probabilistic Forecast Calibration\\
Under Innovation Misspecification
}

\author{Jonathan Ma}

\date{\today}

\begin{document}

\maketitle
\begin{abstract}
    Probabilistic forecasting methods often assume simple innovation distributions, most commonly Gaussian noise, despite real-world systems frequently exhibiting heavy tails, multimodality, heteroskedasticity, and other forms of distributional misspecification. When these assumptions fail, forecast uncertainty may be poorly characterized, leading to degraded calibration and unreliable probabilistic predictions. This thesis investigates whether flexible innovation models can improve probabilistic forecast calibration under innovation misspecification.

    A simulation framework is developed using multivariate vector autoregressive processes with Gaussian, Student-(t), mixture, and heteroskedastic innovation regimes. Forecasting models, including vector autoregressions and probabilistic recurrent neural networks, are fitted to simulated data, and residuals are extracted as empirical proxies for latent innovations. Classical innovation models, including Gaussian, bootstrap, and Student-(t) specifications, are compared against a diffusion-based innovation model trained directly on residual distributions.
    
    Forecast trajectories are generated by propagating sampled innovations through the forecasting models, and probabilistic performance is evaluated using calibration-oriented metrics including empirical coverage, probability integral transform diagnostics, expected calibration error, and proper scoring rules. The results provide insight into the extent to which improved residual distribution learning translates into improved forecast calibration under varying forms of innovation misspecification.
    
    The proposed framework establishes a systematic approach for studying innovation modeling in probabilistic forecasting and provides a foundation for future work on robustness, distributional shift, and calibration-aware innovation learning.
\end{abstract}

\newpage
\tableofcontents
\newpage

\part{Motivation and Background}

\section{Introduction}

Probabilistic forecasting plays a central role in many scientific, economic, and engineering applications. Unlike point forecasting, which provides only a single estimate of a future outcome, probabilistic forecasting seeks to characterize the full uncertainty surrounding future observations. Accurate uncertainty quantification is essential in settings where decisions depend not only on expected outcomes but also on the likelihood of rare or extreme events. Examples include financial risk management, inventory planning, energy systems, healthcare, and environmental forecasting.

A common approach to probabilistic forecasting combines a model for the conditional mean dynamics with an assumption regarding the distribution of future innovations. Classical forecasting models such as autoregressive and vector autoregressive processes frequently assume Gaussian innovations due to their mathematical convenience and analytical tractability. Similarly, many modern machine learning forecasting models employ simple parametric output distributions for uncertainty estimation. While these assumptions often perform adequately in idealized settings, real-world data frequently exhibit departures from Gaussian behavior, including heavy tails, skewness, multimodality, and time-varying volatility.

When the assumed innovation distribution differs substantially from the true data-generating process, forecast uncertainty may be poorly represented even when the conditional mean model is correctly specified. As a result, prediction intervals may be systematically too narrow or too wide, tail probabilities may be misestimated, and probabilistic forecasts may become poorly calibrated. This phenomenon, referred to as innovation misspecification, represents an important but often underexplored source of forecasting error.

Recent advances in generative modeling provide new opportunities for learning complex probability distributions directly from data. In particular, diffusion models have demonstrated remarkable success in approximating high-dimensional distributions across a wide range of applications. Rather than imposing a fixed parametric form, diffusion models learn distributions through an iterative denoising procedure capable of representing complex and highly non-Gaussian behavior. These properties suggest that diffusion models may provide a flexible alternative for innovation modeling within probabilistic forecasting systems.

This thesis investigates whether flexible innovation models can improve probabilistic forecast calibration under innovation misspecification. Rather than replacing the forecasting model itself, the proposed framework focuses on learning the distribution of forecast residuals, which serve as empirical proxies for latent innovations. After fitting forecasting models and extracting residuals, alternative innovation models are trained on the residual distributions and subsequently used to generate future forecast trajectories. This separation between conditional mean modeling and innovation modeling allows the impact of innovation assumptions to be studied directly.

To evaluate this question, a controlled simulation framework is developed using multivariate vector autoregressive processes with multiple innovation regimes, including Gaussian, Student-(t), mixture, and heteroskedastic distributions. Forecasting models are fitted to the simulated data, residuals are extracted, and innovation models are trained using both classical and diffusion-based approaches. Probabilistic forecasts are then generated and evaluated using calibration-oriented metrics such as empirical coverage, expected calibration error, probability integral transform diagnostics, and proper scoring rules.

The central research question of this thesis is:

\begin{quote}
\emph{When do flexible innovation models materially improve probabilistic forecast calibration under innovation misspecification?}
\end{quote}

The primary contribution of this work is the development of a residual-based innovation learning framework that enables direct comparison between classical and diffusion-based innovation models within a common forecasting pipeline. By isolating innovation modeling from conditional mean modeling, the framework provides insight into the relationship between distributional recovery and forecast calibration. The results contribute to a growing literature at the intersection of probabilistic forecasting, generative modeling, and uncertainty quantification.

The remainder of the thesis is organized as follows. Chapter 2 reviews the relevant literature on innovation modeling, forecast calibration, diffusion models, and robustness under distributional shift. The empirical framework and simulation environments are introduced in Part II. Classical and diffusion-based innovation models are developed and evaluated in Parts III and IV. Calibration and robustness experiments are presented in Part V. The thesis concludes with a discussion of findings, limitations, and directions for future research.
\section{Literature Review}

\textbf{\href{https://link.springer.com/book/10.1007/978-3-540-27752-1}{Lütkepohl (2005)}} provides the classical forecasting foundation used throughout this thesis. Vector autoregressions (VARs) model future observations using lagged values of a multivariate time series while representing uncertainty through an innovation process. Under the moving-average representation, forecast uncertainty propagates entirely through future innovations. Consequently, assumptions regarding the innovation distribution directly determine the quality of probabilistic forecasts and prediction intervals. This observation motivates the central question of the thesis: whether more flexible innovation models can improve forecast calibration when classical Gaussian assumptions are violated.

\textbf{\href{https://arxiv.org/abs/1704.04110}{Salinas et al. (2020)}} introduced DeepAR, a deep-learning approach to probabilistic forecasting that learns a shared recurrent neural network across multiple time series. Forecast distributions are generated through autoregressive Monte Carlo sampling rather than closed-form prediction intervals. The paper demonstrates that forecasting can be viewed as a distribution-learning problem rather than a point-prediction problem. This thesis adopts a similar probabilistic forecasting perspective while focusing specifically on innovation uncertainty.

\textbf{\href{https://arxiv.org/abs/2006.11239}{Ho et al. (2020)}} introduced Denoising Diffusion Probabilistic Models (DDPMs), a flexible generative framework capable of learning highly complex distributions through iterative denoising. Rather than imposing a parametric density, diffusion models learn a reverse stochastic process that transforms noise into samples from a target distribution. Their ability to represent heavy-tailed, multimodal, and non-Gaussian distributions motivates their use as innovation models in this thesis.

\textbf{\href{https://arxiv.org/abs/2011.13456}{Song et al. (2021)}} unified diffusion models and score-based generative modeling through stochastic differential equations. This work established diffusion learning as score estimation and provided the theoretical foundation for modern diffusion methods. The score-based interpretation offers a mathematical framework for treating innovation learning as a generative modeling problem and motivates the diffusion methodology employed throughout the empirical experiments.

\textbf{\href{https://arxiv.org/abs/2101.12072}{Rasul et al. (2021)}} extended diffusion models to probabilistic time-series forecasting through TimeGrad. Their work demonstrated that diffusion models can successfully represent complex forecast distributions and outperform simpler parametric uncertainty models. TimeGrad serves as an important methodological reference because it established diffusion models as viable forecasting tools. However, unlike TimeGrad, the present thesis does not model future observations directly. Instead, diffusion models are applied to residual innovations obtained after fitting a forecasting model.

\textbf{\href{https://fitelson.org/seminar/dawid.pdf}{Dawid (1982)}} provided one of the foundational definitions of probabilistic calibration. Calibration requires that predicted probabilities agree with observed frequencies over repeated forecasting exercises. This concept serves as the primary evaluation target of the thesis. A flexible innovation model is only useful if it improves the calibration of the resulting forecast distributions.

\textbf{\href{https://arxiv.org/abs/1706.04599}{Guo et al. (2017)}} demonstrated that highly accurate machine learning models can nevertheless produce poorly calibrated uncertainty estimates. The paper introduced Expected Calibration Error (ECE), which remains one of the most widely used calibration metrics. Their findings motivate evaluating probabilistic forecasts using calibration measures rather than predictive accuracy alone.

\textbf{\href{https://sites.stat.washington.edu/raftery/Research/PDF/Gneiting2007jasa.pdf}{Gneiting and Raftery (2007)}} established proper scoring rules as the standard framework for evaluating probabilistic forecasts. They introduced metrics such as the Continuous Ranked Probability Score (CRPS) and Energy Score while emphasizing the tradeoff between calibration and sharpness. These metrics provide the primary evaluation framework used throughout the empirical analysis.

\textbf{\href{https://arxiv.org/abs/1912.08335}{Masegosa (2019)}} studies predictive inference under model misspecification and argues that likelihood-based learning objectives do not necessarily yield optimal predictive uncertainty. The paper highlights the distinction between accurately fitting a model and producing reliable predictive distributions. This perspective motivates the thesis focus on innovation misspecification. Even when forecasting models are correctly specified, incorrect innovation assumptions may lead to poorly calibrated uncertainty estimates. The empirical experiments therefore investigate whether flexible diffusion-based innovation models can mitigate calibration degradation under misspecification.

Collectively, these works motivate the central thesis of this project. Classical forecasting methods derive predictive uncertainty from assumptions on the innovation process, while modern diffusion models provide a flexible mechanism for learning complex innovation distributions directly from data. At the same time, calibration literature emphasizes that accurate probabilistic forecasting requires more than expressive generative models; predictive distributions must also align with observed frequencies. This thesis therefore investigates whether diffusion-based innovation models can improve probabilistic forecast calibration under innovation misspecification when compared to traditional parametric alternatives.

\part{Empirical Experiments}

\section{Simulation Framework}

\section{Innovation Modeling}

\section{Forecast Generation}

\part{Calibration and Robustness}

\section{Calibration and Evaluation}


\section{Robustness Testing}

\part{Conclusion and Discussion}

\section{Discussion}

\section{Conclusion}

\section{Future Work}

\part{Appendix}

\appendix

\section{Tables}

\section{Figures}

\end{document}